\section{Fondamenti della modellazione orientata agli oggetti}
I fondamenti della modellazione orientata agli oggetti si basano sul concetto di \textbf{oggetto}.

Ogni oggetto nel sistema rappresenta un'istanza di una specifica classe, la quale definisce preventivamente l'insieme comune di caratteristiche condivise da tutte le sue istanze

\subsection{Elementi di un oggetto}
Tutti gli oggetti sono definiti da 3 elementi imprescindibili
\subsubsection{Identità}
Ogni oggetto possiede un proprio identificativo univoco che lo distingue in modo inequivocabilmente dagli altri
\subsubsection{Stato (dati)}
È determinato dai valori effettivi dei dati memorizzati negli attributi dell'oggetto in un preciso istante temporale
\subsubsection{Comportamento (le operazioni)}
Rappresenta l'insieme delle azioni e delle operazione che l'ggetto è in grado di eseguire

\subsection{Incapsulamento}
Qeusta regola stabilisce che i dati devono essere nascosti all'interno dell'oggetto

Di conseguenza, l'unico modo per avvedervi o modificarli è attraverso l'invocazione delle operazioni pubbliche dell'oggetto stesso

\subsection{Nessaggi}
Gli oggetti non operano in modo isolato, ma collaborano per eseguire le funzioni del sistema scambiandosi messaggi.

Quando un oggetto invia un messaggio a un altro oggetto, tale messaggio causa l'invocazione della corrispondente operazione nel destinatario
\subsubsection{Tipologie di messaggi}
\begin{figure}[H]
    \centering
    \includegraphics[width=\linewidth]{Schermata_20260428_162015.png}
\end{figure}


\section{Diagrammi di Interazione}
Il termine diagramma di interazione è una generalizzazione di 2 tipi più specifici di diagrammi di UML

\begin{itemize}
    \item \textbf{Diagrammi di Sequenza}

    Mostrano il le interazioni in una specie di formato a steccato, in cui gli oggetti che partecipano all'interazione sono mostrati in alto, uno a fianco all'altro

    \begin{figure}[H]
        \centering
        \includegraphics[width=0.6\linewidth]{Schermata_20260428_161345.png}
    \end{figure}
    \item \textbf{Diagrammi di comunicazione}
    Mostrano le interazioni tra gli oggetti in un formato a grafo o rete, in cui gli oggetti possono essere posizionati ovunque nel diagramma
    \begin{figure}[H]
        \centering
        \includegraphics[scale=0.8]{Schermata_20260428_161707.png}
    \end{figure}
\end{itemize}
\subsection{Vantaggi e Svantaggi}
\begin{tabularx}{\linewidth}{|l | X | X|}
    \hline Tipo & Punti di forza & Punti di debolezza \\ \hline
    Diagramma di sequenza & Mostra chiaramente la sequenza o l'ordinamento temporale dei messaggi. 

    Opzioni di notazione numerose e dettagliate & Costringe a estendersi verso destra quando si aggiungo nuovi oggetti 

    Consuma lo spazio orizzontale \\ \hline
    Diagramma di comunicazione & Economizza lo spazio

    flessibilità nell'aggiunta di nuovi oggetti nelle 2 dimensioni & Più difficile vedere la sequenza dei messaggi.

    Meno opzioni di notazione.\\ \hline
\end{tabularx}
